\chapter{Experiment definition with AEDL}\label{ch:aedl-reference}

This chapter lists all AEDL directives. Defaults are given in brackets.

\section{Source and integrated luminosity}\label{sec:aedl-flux}

\begin{center}\small
\begin{tabularx}{\linewidth}{lX}
\toprule
\multicolumn{2}{l}{\code{nuflux(\#name)< ... >} (or the older \code{flux}, see below)} \\
\midrule
\code{@flux\_file = "file"} & flux table (see below) \\
\code{@builtin = n} & built-in flux: 1 $\mu^+$ decay, 2 $\mu^-$ decay (neutrino factory),
  3 $\beta^+$ decay, 4 $\beta^-$ decay (beta beam) \\
\code{@time} & running time [1] \\
\code{@power} & source power [1], e.g.\ MW \\
\code{@stored\_muons}, \code{@stored\_ions} & useful parent decays per year [1] \\
\code{@norm} & normalisation of the flux table [1] \\
\code{@parent\_energy} & muon energy in GeV (built-in 1, 2) \\
\code{@gamma}, \code{@end\_point} & Lorentz factor of the ions and endpoint of the
  decay in GeV (built-in 3, 4) \\
\midrule
\code{\$target\_mass} & fiducial detector mass, e.g.\ kt \\
\bottomrule
\end{tabularx}
\end{center}

\paragraph{Flux files} have 501 lines of seven columns, with equidistant
energies:
\[
E\ [\mathrm{GeV}]\quad \phi_{\nue}\quad\phi_{\numu}\quad\phi_{\nu_\tau}\quad
\phi_{\bar\nu_e}\quad\phi_{\bar\nu_\mu}\quad\phi_{\bar\nu_\tau}
\]
Lines starting with \code{\#} are comments. The flux is interpolated linearly
and is zero outside the table. Unused flavours must be given as zeros.

\paragraph{Normalisation.} The flux in Eq.~\eqref{eq:flux} is
$\mathcal N/L^2\times$\code{@norm}$\times$\code{@power}$\times$\code{@stored\_muons}$\times$\code{@time}$\times\phi(E)$,
with $\mathcal N=1$ for \code{nuflux} tables. Together with the cross sections
in $10^{-38}\,\mathrm{cm^2}$ per nucleon and the target mass $M$, the number
of events per GeV is
\begin{equation}
\frac{dN}{dE} = \frac{\texttt{@norm}\times\texttt{@power}\times\texttt{@time}\times M}{(L/\mathrm{km})^2}\,
  \phi(E)\,\frac{\sigma(E)}{10^{-38}\,\mathrm{cm}^2}\,P(E).
\end{equation}
If the table gives the number of neutrinos per GeV and per area $A$ at a
distance $L_0$ for an integrated luminosity $\mathcal L$ (in the units
$\mathcal L_u$ chosen for \code{@power}$\times$\code{@time}), the normalisation is
\begin{equation}
\texttt{@norm} = \frac{\mathrm{cm}^2}{A}\left(\frac{L_0}{\mathrm{km}}\right)^2
  \frac{N_\text{target}}{M_u}\times10^{-38}\times\frac{\mathcal L_u}{\mathcal L},
\end{equation}
with $N_\text{target}/M_u$ the number of target nucleons per unit of the
target mass ($6.022\times10^{32}$ per kt). The older environment \code{flux}
multiplies by $\mathcal N = 5.1989$ in addition; replace \code{flux} by
\code{nuflux} and multiply \code{@norm} by $5.1989$ to convert. Built-in fluxes
have their own normalisation and need no \code{@norm}.

\begin{figure}[t]
\centering
\begin{tikzpicture}
\begin{axis}[manual,xlabel={$E$ [GeV]},ylabel={$\phi$ at 1~km [a.u.]},xmin=0,xmax=50,ymin=0,
  title={\small neutrino factory, $\mu^+$, 50 GeV},legend pos=north west]
\addplot table[x=E,y=numubar]{figures/data/nf_flux.dat};
\addplot table[x=E,y=nue]{figures/data/nf_flux.dat};
\legend{$\bar\nu_\mu$,$\nue$}
\end{axis}
\end{tikzpicture}\hfill
\begin{tikzpicture}
\begin{semilogxaxis}[manual,xlabel={$E$ [GeV]},ylabel={$\sigma/E$ [$10^{-38}\,\mathrm{cm^2/GeV}$]},
  xmin=0.1,xmax=100,ymin=0,title={\small charged current, \file{XCC.dat}},legend pos=north west]
\addplot table[x=E,y=numu]{figures/data/xsec.dat};
\addplot table[x=E,y=numubar]{figures/data/xsec.dat};
\legend{$\numu$,$\bar\nu_\mu$}
\end{semilogxaxis}
\end{tikzpicture}
\caption{Left: the built-in neutrino factory flux (\code{@builtin = 1}) of
\file{NFstandard.glb}, \code{flux(nf, E; flavour=:μ, antineutrino=true, L=1)}.
Right: the charged-current cross sections per nucleon divided by energy,
\code{cross\_section(nf, E; flavour=:μ) / E}.}
\label{fig:flux}
\end{figure}

\paragraph{Built-in fluxes.} The neutrino factory flux (\code{@builtin = 1, 2})
is that of muons of energy $E_\mu$ (\code{@parent\_energy}) decaying at rest in
the beam direction, with $x = 2E/(E_\mu(1+\beta))$,
\begin{equation}
\phi_{\bar\nu_\mu}\ (\phi_{\numu}) \propto \frac{\gamma(1+\beta)}{\pi m_\mu}\,x^2(3-2x),\qquad
\phi_{\nue}\ (\phi_{\bar\nu_e}) \propto \frac{\gamma(1+\beta)}{\pi m_\mu}\,6x^2(1-x),
\end{equation}
for $\mu^+$ ($\mu^-$) decays, up to the endpoint $x=1$. The beta beam flux
(\code{@builtin = 3, 4}) is the spectrum of nuclear beta decay with endpoint
$E_0$ (\code{@end\_point}), boosted by $\gamma$ (Figure~\ref{fig:flux}):
\begin{equation}
\phi(E) = \frac{\gamma^2}{\pi\,g(y_e)\,E}\,x^3(1-x)\sqrt{(1-x)^2-y_e^2},\qquad
x = \frac{E}{2\gamma(E_0+m_e)},\quad y_e = \frac{m_e}{E_0+m_e},
\end{equation}
with the phase space integral
$g(y)=\frac1{60}\bigl[\sqrt{1-y^2}(2-9y^2-8y^4)+15y^4\ln\frac{1+\sqrt{1-y^2}}{y}\bigr]$
(normalised to a distance of 1~km).

\begin{cdiff}
C GLoBES uses $\pi=3.1415$, $\log_{10}$ instead of $\ln$ in $g(y)$, an old value
of Avogadro's number, and rounded muon and electron masses. Its flux
interpolation divides the energy range of the table into 501 instead of 500
steps, which evaluates every flux at $1.002$ times the requested energy, and
the built-in fluxes are tabulated only up to $500/501$ of their endpoint
(Section~\ref{sec:cdiff-flux}).
\end{cdiff}

\section{Baseline and matter density profile}\label{sec:aedl-baseline}

\begin{center}\small
\begin{tabularx}{\linewidth}{lX}
\toprule
\code{\$baseline} & baseline in km \\
\code{\$profiletype = 1} & a single layer with the average density along \code{\$baseline} \\
\code{\$profiletype = 2} & \code{\$densitysteps} layers of equal length with the Earth's density \\
\code{\$profiletype = 3} & layers \code{\$lengthtab = \{...\}} (km) with densities
  \code{\$densitytab = \{...\}} ($\mathrm{g/cm^3}$); the baseline is their total length \\
\bottomrule
\end{tabularx}
\end{center}
A constant density $\rho$ is \code{\$profiletype = 3},
\code{\$densitytab = \{ρ\}}, \code{\$lengthtab = \{L\}}. See
Section~\ref{sec:profiles}.

\section{Cross sections}\label{sec:aedl-cross}

\code{cross(\#name)< @cross\_file = "file" >} reads a table of 1001 lines with
equidistant $\log_{10}(E/\mathrm{GeV})$:
\[
\log_{10}E\quad \hat\sigma_{\nue}\quad\hat\sigma_{\numu}\quad\hat\sigma_{\nu_\tau}\quad
\hat\sigma_{\bar\nu_e}\quad\hat\sigma_{\bar\nu_\mu}\quad\hat\sigma_{\bar\nu_\tau},
\qquad \hat\sigma = \frac{\sigma(E)}{E}\ \bigl[10^{-38}\,\mathrm{cm^2/GeV}\bigr].
\]
The table is interpolated linearly in $\log_{10}E$ and is zero outside. The
files shipped with GLoBES, \file{XCC.dat} (charged current), \file{XNC.dat}
(neutral current) and \file{XQE.dat} (quasi-elastic), cover $0.1$ to $1000\GeV$
and are per nucleon of an isoscalar target.

\section{Channels}\label{sec:aedl-channel}

\begin{aedl}
channel(#name)<
   @channel = #flux : +|- : initial : final : #cross : #energy
   @pre_smearing_efficiencies  = {...}    // one per sampling point [1]
   @pre_smearing_background    = {...}    // one per sampling point [0]
   @post_smearing_efficiencies = {...}    // one per bin [1]
   @post_smearing_background   = {...}    // one per bin [0]
>
\end{aedl}
\code{+} and \code{-} are neutrinos and antineutrinos of the flux. Flavours are
\code{e}, \code{m}, \code{t} (also \code{electron}, \code{muon}, \code{tau});
\code{NOSC\_e} etc.\ give a channel without oscillations ($P=1$ for equal
flavours), for neutral currents and intrinsic backgrounds. The efficiencies and
backgrounds are $\epsilon^\text{pre}$, $B^\text{pre}$, $\epsilon^\text{post}$
and $B^\text{post}$ of Eqs.~\eqref{eq:presmear} and~\eqref{eq:postsmear}.

\section{Energy resolution}\label{sec:aedl-energy}

\paragraph{Binning.}
\begin{center}\small
\begin{tabularx}{\linewidth}{lX}
\toprule
\code{\$emin}, \code{\$emax} & range of reconstructed energies in GeV \\
\code{\$bins} & number of bins of equal width \\
\code{\$binsize = \{...\}} & bin widths instead, for unequal bins (they must add up to the range) \\
\code{\$sampling\_min}, \code{\$sampling\_max} & range of true energies [\code{\$emin}, \code{\$emax}] \\
\code{\$sampling\_points} & number of sampling points [\code{\$bins}] \\
\code{\$sampling\_stepsize = \{...\}} & widths of the sampling steps instead \\
\bottomrule
\end{tabularx}
\end{center}

\paragraph{Resolution functions.}
\begin{aedl}
energy(#name)<
   @type = 1                       // 2: reduce the width by the filter
   @sigma_function = #standard     // or #inverse_beta
   @sigma_e = {0.15, 0, 0}         // σ(E) = 0.15 E
>
\end{aedl}
The smearing matrix is computed from Eq.~\eqref{eq:smear} with the
resolution functions of Section~\ref{sec:smearing}.

\paragraph{Smearing matrices.} The matrix $K_{ij}$ can also be given row by
row, one row per bin $i$, each with the first and last sampling point (counted
from 0) of its non-zero entries followed by the entries:
\begin{aedl}
energy(#migration)<
   @energy = {0, 2, 0.6, 0.3, 0.1} :
             {0, 3, 0.2, 0.5, 0.2, 0.1} :
             ...
             {17, 19, 0.1, 0.3, 0.6} ;
>
\end{aedl}
Matrices from detector simulations (for example the beta beam files
\file{*EffMig350.dat}) are included this way.

\paragraph{Filter.}\label{sec:aedl-filter} \code{\$filter\_state = 1} and
\code{\$filter\_value = }$\sigma_f$ (GeV) switch on the low-pass filter of
Section~\ref{sec:filter}.

\section{Rules and systematics}\label{sec:aedl-rule}

\begin{aedl}
rule(#name)<
   @signal          = 0.5 @ #channel_1
   @background      = 0.001 @ #channel_2 : 0.005 @ #channel_3
   @signalerror     = 0.05 : 0.0001        // σ1 : σ2
   @backgrounderror = 0.1 : 0.0001         // σ3 : σ4
   @backgroundcenter = 1.0 : 0.0           // β0 : τ0 (better not used)
   @sys_on_function  = "chiSpectrumTilt"
   @sys_off_function = "chiNoSysSpectrum"
   @energy_window    = 0.5 : 3.0            // GeV
>
\end{aedl}
\code{@signal} and \code{@background} are the coefficients $a_k$ and $b_k$ of
Eq.~\eqref{eq:rule}. The errors are those of the built-in $\chi^2$ functions
(Section~\ref{sec:builtin-chi2}). Functions with other numbers of parameters
(user-defined ones) take their errors as lists, \code{@sys\_on\_errors =
\{0.02, 0.05\}} and \code{@sys\_off\_errors}, which also override
\code{@signalerror} and \code{@backgrounderror} for built-in functions. The
energy window restricts the bins used in the $\chi^2$ of the rule
[\code{\$emin} : \code{\$emax}]; with an energy calibration error, it should
leave a margin of a few times the error at both ends. GLoBES 2 files may use
\code{@errordim\_sys\_on} and \code{@errordim\_sys\_off}
(Table~\ref{tab:chi2functions}).

\section{Testing AEDL files}

Load the file and look at what it defines:
\begin{julia}
e = Experiment("MyExperiment.glb")       # parse errors are thrown with file and line
display(e)                               # baseline, mass, bins, rules
channel_names(e), rule_names(e)
simulate!(truth)
event_rates(e; rule=1)                   # signal and background per bin
[sum(channel_rates(e, c)) for c in eachindex(channel_names(e))]
\end{julia}
Compare the total event numbers with those published for the experiment; most
problems are a wrong flux normalisation (\code{@norm}) or efficiencies.
